\chapter{Fachdidaktischer und theoretischer Rahmen}\label{chap:fachdidaktischer-rahmen}
\todo[color=yellow]{Theoretische Grundlagen und Fachdidaktik ausarbeiten.}

\section{Curriculare Verortung}\label{sec:curriculare-verortung}
\todo[color=yellow]{Lehrplanbezug und Kompetenzen zuordnen.}

\section{Interdisziplinarität}\label{sec:interdisziplinaritaet}
\todo[color=yellow]{Fächerübergreifende Bezüge herausarbeiten.}

\subsection{Bezüge zur Physik}\label{subsec:bezuege-physik}
\todo[color=yellow]{Strom, Spannung und Leistung einordnen.}

\subsection{Bezüge zur Mathematik}\label{subsec:bezuege-mathematik}
\todo[color=yellow]{Logik und Messdaten mathematisch verknüpfen.}

\section{Physical Computing und Tangible Interfaces}\label{sec:physical-computing}
\todo[color=yellow]{Tangible Interfaces lerntheoretisch begründen.}

\section{Prinzipien didaktischer Reduktion}\label{sec:didaktische-reduktion}
\todo[color=yellow]{Low Floor und High Ceiling erklären.}
